\documentclass[a5paper,10pt]{article}

% https://github.com/InDevRus/filippov-solutions
% Нумерация страниц
\pagenumbering{gobble}

% Настройка полей
\usepackage{geometry}
\geometry{left=1cm}
\geometry{right=1cm}
\geometry{top=1cm}
\geometry{bottom=1.5cm}

% Вставка таблицы
\usepackage{array}
\usepackage{tabularx}

% Инструменты выравнивания формул
\usepackage{amsmath}
\usepackage{mathtools}

% Диакритика в математике
\usepackage{amsfonts}

% Вычисления размеров на ходу
\usepackage{calc}

% Удобные записи производных
\usepackage[italic=true]{derivative}

% Настройки сносок
\usepackage{enotez}

% Длинные знаки векторов
\usepackage{anyfontsize}
\usepackage[b]{esvect}

% Вставка картинок
\usepackage{graphicx}

% Вставка красной строки
\usepackage{indentfirst}
\setlength{\parindent}{1em}

% Условия в командах
\usepackage{xifthen}

% Луа-скрипты
\usepackage{luacode}

% Настройка команд отдельно в математическом и текстовом режимах
\usepackage{mathcommand}

% Для повторения LaTeX кода
\usepackage{multido}

% Конфигурация отступов формул
\delimitershortfall-1sp
\usepackage{mleftright}
\mleftright

% Растягивание объектов
\usepackage{relsize}

% Обтекание текстом
\usepackage{wrapfig}
\usepackage{subcaption}
\usepackage{floatflt}

% Поддержка ввода кириллицы
\usepackage[english, russian]{babel}

% Настройка корневой позиции для компилятора
\usepackage{import}
\providecommand*{\ProjectRootPath}{..}

\import{\ProjectRootPath/preambles/macros/}{theorems}

\import{\ProjectRootPath/preambles/fonts}{font_setup}

% Знак градуса
\usepackage{siunitx}

\import{\ProjectRootPath/preambles/macros/}{operators}
\import{\ProjectRootPath/preambles/macros/}{redefinitions}

\import{\ProjectRootPath/preambles/fonts}{letters_setup}
\import{\ProjectRootPath/preambles/fonts/adjustments}{custom_superscripts}

\import{\ProjectRootPath/preambles/custom}{formatting}
\import{\ProjectRootPath/preambles/custom}{frames}
\import{\ProjectRootPath/preambles/custom}{list_setup}

% TODO: перевести команды на современный стиль объявления

\begin{document}

$y' = \dfrac {1 + \cos 2y} {x^2}$.
Функции $y\br{x}$ такие, что $\cos 2y = -1$,
являются решениями, и такими функциями
будут $y\br{x} = -\dfrac {\pi} {2} + \pi n$, $n$ -- целое.
$\dfrac {y'} {1 + \cos 2y} = \dfrac {1} {x^2}$.
$\tint \dfrac {dx} {x^2} = -\dfrac {1} {x} + C$.
$$\tint \dfrac {y'} {1 + \cos\br{2y\br{x}}} dx
= \tint \dfrac {dy} {1 + \cos 2y}
= \tint \dfrac {dy} {2\cos^2 y}
= \dfrac {1} {2} \tg y + C .$$

Получаем $\tg y = -\dfrac {2} {x} + C$. Общее решение имеет вид $y\br{x} = \arctg \br{- \dfrac {2} {x} + C} + \pi n$,
а также дополнительные $y\br{x} = -\dfrac {\pi} {2} + \pi n$.

Теперь найдём частные решения, для которых
$\tlim{\mathsmaller{x \to +\infty}} y\br{x} = \dfrac {9\pi} {4}$.
\begin{itemize}
	\item Для общего решения имеем $\tlim{\mathsmaller{x \to +\infty}} y\br{x} = \arctg C + \pi n = \dfrac {9\pi} {4}$. Найдём сначала $n$.
	$\dfrac {9\pi} {4} - \dfrac {\pi} {2} < \dfrac {9\pi} {4} - \arctg C < \dfrac {9\pi} {4} + \dfrac {\pi} {2}$, так что $\dfrac {7\pi} {4} < \pi n < \dfrac {11\pi} {4}$. \vspace{1mm}
	Единственное подходящее $n$ -- это $2$. Отсюда $\arctg C = \dfrac {9\pi} {4} - 2\pi = \dfrac {\pi} {4}$ и $C = 1$.
	
	Получаем частное решение $y_{1}\br{x} = \arctg\br{1 - \dfrac {2} {x}} + 2\pi$.
	
	\item Для дополнительного решения имеем $\tlim{\mathsmaller{x \to +\infty}} = -\dfrac {\pi} {2} + \pi n = \dfrac {9\pi} {4}$. Таких целых $n$ нет, значит среди дополнительных решений подходящих не найдётся.
\end{itemize}

\end{document}
